% Figure: a rotary regenerator (thermal wheel) that stores heat in a slowly
% rotating porous matrix, one half swept by the hot stream and the other half by
% the cold stream, so each matrix element alternately absorbs heat from the hot
% side and releases it to the cold side as the wheel turns.
\begin{figure}[htbp]
\centering
\begin{tikzpicture}[>=stealth]
% the wheel
\draw[thick] (0,0) circle (2.2);
\draw[thick] (0,0) circle (0.35);
% split the wheel into hot and cold halves along a vertical diameter
\draw[thick, dashed] (0,-2.2) -- (0,2.2);
% hot half shaded on the left, cold half on the right
\begin{scope}
  \clip (0,0) circle (2.2);
  \fill[engred!18] (-2.2,-2.2) rectangle (0,2.2);
  \fill[engblue!18] (0,-2.2) rectangle (2.2,2.2);
\end{scope}
\draw[thick] (0,0) circle (2.2);
\draw[thick] (0,0) circle (0.35);
% matrix hatching suggested by spokes
\foreach \a in {30,60,120,150,210,240,300,330} {
  \draw[enggray!60] (\a:0.35) -- (\a:2.2);
}
% rotation arrow
\draw[->, thick, engamber] (150:1.3) arc (150:30:1.3);
\node[engamber, font=\scriptsize] at (90:1.55) {rotation};
% hot stream through left half
\draw[engred, ->, very thick] (-3.4,1.1) -- (-2.4,1.1);
\draw[engred, ->, very thick] (-1.6,1.1) -- (-0.6,1.1);
\node[engred, font=\scriptsize, anchor=south] at (-3.0,1.2) {hot gas in};
\draw[engred, ->, thick] (-2.4,-1.1) -- (-3.4,-1.1);
\node[engred, font=\scriptsize, anchor=north] at (-3.0,-1.2) {cooled out};
% cold stream through right half
\draw[engblue, ->, very thick] (3.4,-1.1) -- (2.4,-1.1);
\draw[engblue, ->, very thick] (1.6,-1.1) -- (0.6,-1.1);
\node[engblue, font=\scriptsize, anchor=north] at (3.0,-1.2) {cold air in};
\draw[engblue, ->, thick] (2.4,1.1) -- (3.4,1.1);
\node[engblue, font=\scriptsize, anchor=south] at (3.0,1.2) {heated out};
% labels for the two halves
\node[engred!70!black, font=\scriptsize] at (-1.1,0.0) {matrix absorbs};
\node[engblue!70!black, font=\scriptsize] at (1.1,0.0) {matrix releases};
\end{tikzpicture}
\caption[Rotary regenerator (thermal wheel)]{A rotary regenerator stores heat in
a slowly turning porous matrix rather than passing it through a wall. The hot
stream sweeps one half of the wheel and warms the matrix it passes; the wheel
carries that warmed matrix into the cold half, where the cold stream sweeps the
stored heat back out; and the cooled matrix returns to the hot side to repeat the
cycle. Each element of the matrix is therefore a small heat store that is charged
and discharged once per revolution, so the device is a regenerator, storing heat
in time, rather than a recuperator that transfers it in space through a fixed
wall. The trade the geometry carries is carry-over: a thin sliver of one stream
is trapped in the matrix pores and dragged across the seal into the other stream
at each pass, which rules the wheel out where the two streams must not mix and
makes it the standard choice for building ventilation where a little air
carry-over is harmless.}
\label{fig:vol04ch06:rotary-regenerator}
\end{figure}
